\documentclass[11pt]{article}

\usepackage[margin=1in]{geometry}
\usepackage{amsmath,amssymb,amsthm,mathtools}
\usepackage{enumitem}

\newtheorem{definition}{Definition}[section]
\newtheorem{lemma}[definition]{Lemma}
\newtheorem{proposition}[definition]{Proposition}
\newtheorem{theorem}[definition]{Theorem}
\newtheorem{remark}[definition]{Remark}

\newcommand{\R}{R}
\newcommand{\Sn}{\mathcal S}
\newcommand{\pow}[1]{R^{#1}}
\newcommand{\Graph}{\operatorname{Graph}}
\newcommand{\Def}{\operatorname{Def}}
\newcommand{\pr}{\operatorname{pr}}
\newcommand{\dom}{\operatorname{dom}}
\newcommand{\LC}{\operatorname{LC}}
\newcommand{\Cont}{\operatorname{Cont}}
\newcommand{\LInc}{\operatorname{LInc}}
\newcommand{\LDec}{\operatorname{LDec}}
\newcommand{\val}{\operatorname{val}}
\newcommand{\nhd}{\operatorname{nhd}}
\newcommand{\core}{\operatorname{core}}

\title{Human-language version of the Lean formalisation\\
O-minimal structures, definable functions, and local properties}
\author{}
\date{}

\begin{document}
\maketitle

\section{Ambient ordered set and powers}

\begin{definition}[Dense linear order without endpoints]
A \emph{dense linear order without endpoints} is a type/set $R$ equipped with a binary relation $<$ such that:
\begin{enumerate}[label=(\roman*)]
  \item $\neg(x<x)$ for every $x\in R$;
  \item if $x<y$ and $y<z$, then $x<z$;
  \item for all $x,y\in R$, exactly one of $x<y$, $x=y$, or $y<x$ holds;
  \item if $x<y$, then there is $z\in R$ with $x<z<y$;
  \item for every $x\in R$, there are $y,z\in R$ with $y<x<z$.
\end{enumerate}
This is the mathematical content of the Lean structure
\texttt{DenseLinearOrderNoEndpoints}.
\end{definition}

For a natural number $n$, we write $R^n$ for the set of $n$-tuples of elements of $R$.
In Lean this is implemented as functions
\[
  \operatorname{Fin}(n)\to R,
\]
because a tuple of length $n$ is equivalently a function whose domain is the finite set of $n$ indices.

If $x\in R^n$ and $y\in R^m$, their concatenation is written
\[
  x\mathbin{{}^\frown}y\in R^{n+m}.
\]
If $z\in R^{n+m}$, then $z_L\in R^n$ and $z_R\in R^m$ denote the left and right coordinate blocks. Thus
\[
  (x\mathbin{{}^\frown}y)_L=x,
  \qquad
  (x\mathbin{{}^\frown}y)_R=y.
\]
These are the human versions of the Lean definitions
\texttt{Power.append}, \texttt{Power.left}, and \texttt{Power.right}.

For subsets we use the ordinary operations
\[
  \emptyset,\qquad A\cup B,\qquad A\cap B,
  \qquad R^n\setminus A.
\]
In Lean these are implemented directly as predicates:
\[
  \emptyset(x)\equiv \bot,
  \quad
  (A\cup B)(x)\equiv A(x)\lor B(x),
  \quad
  (A\cap B)(x)\equiv A(x)\land B(x),
  \quad
  A^c(x)\equiv \neg A(x).
\]

\section{Intervals and one-dimensional definable sets}

To state the o-minimality axiom, the Lean file introduces formal endpoints
\[
  -\infty,\qquad a\in R,
  \qquad +\infty.
\]
The order on endpoints is the expected one: $-\infty$ is below all finite points, $+\infty$ is above all finite points, and finite endpoints are compared using the given order on $R$.

For endpoints $a,b\in R\cup\{-\infty,+\infty\}$, the corresponding open interval is
\[
  (a,b)=\{x\in R:a<x<b\},
\]
with the obvious interpretation at infinite endpoints.
A point-set is a singleton
\[
  \{a\}=\{x\in R:x=a\}.
\]

\begin{definition}[Finite unions of points and intervals]
A subset $A\subseteq R$ is a finite union of points and intervals if it belongs to the smallest class containing:
\begin{enumerate}[label=(\roman*)]
  \item the empty set;
  \item every singleton $\{a\}$;
  \item every open interval $(a,b)$ with possibly infinite endpoints;
  \item the union of any two members of the class.
\end{enumerate}
This is the inductive predicate \texttt{FiniteUnionOfPointsAndIntervals} in the Lean file.
\end{definition}

\section{O-minimal structures}

\begin{definition}[O-minimal structure]
Let $(R,<)$ be a dense linear order without endpoints. An \emph{o-minimal structure} on $R$ consists of a family
\[
  \mathcal S_n\subseteq \mathcal P(R^n),
  \qquad n=0,1,2,\ldots,
\]
satisfying the following closure properties.
\begin{enumerate}[label=(\roman*)]
  \item $\emptyset\in \mathcal S_n$ for every $n$.
  \item If $A,B\in\mathcal S_n$, then $A\cup B\in\mathcal S_n$.
  \item If $A,B\in\mathcal S_n$, then $A\cap B\in\mathcal S_n$.
  \item If $A\in\mathcal S_n$, then $R^n\setminus A\in\mathcal S_n$.
  \item For $i<j<n$, the diagonal
  \[
    \{x\in R^n:x_i=x_j\}
  \]
  belongs to $\mathcal S_n$.
  \item If $A\in\mathcal S_n$ and $B\in\mathcal S_m$, then
  \[
    A\times B
    =\{z\in R^{n+m}:z_L\in A\text{ and }z_R\in B\}
  \]
  belongs to $\mathcal S_{n+m}$.
  \item If $A\in\mathcal S_{n+1}$, then each one-coordinate projection of $A$ belongs to $\mathcal S_n$.
  \item Reindexing coordinates preserves definability: if $A\in\mathcal S_n$ and
  $\sigma:\{0,\ldots,n-1\}\to\{0,\ldots,m-1\}$, then
  \[
    \{x\in R^m:(x_{\sigma(0)},\ldots,x_{\sigma(n-1)})\in A\}
    \in\mathcal S_m.
  \]
  \item Block existential projection preserves definability: if $A\in\mathcal S_{n+m}$, then
  \[
    \{x\in R^n:\exists y\in R^m,\,x\mathbin{{}^\frown}y\in A\}
    \in\mathcal S_n.
  \]
  \item The order relation
  \[
    \{(x,y)\in R^2:x<y\}
  \]
  belongs to $\mathcal S_2$.
  \item O-minimality: for every $A\subseteq R$,
  \[
    A\in\mathcal S_1
    \quad\Longleftrightarrow\quad
    A\text{ is a finite union of points and intervals.}
  \]
\end{enumerate}
This is the Lean structure \texttt{OMinimalStructure}.
\end{definition}

\begin{remark}
Some closure properties above, especially reindexing and block existential projection, are often derived from more primitive axioms. They are included explicitly in the Lean file to avoid a large amount of finite-coordinate bookkeeping. All subsequent arguments use only the listed operations.
\end{remark}

A set $A\subseteq R^n$ is called \emph{definable} if $A\in\mathcal S_n$.

\section{Definable functions and composition}

\begin{definition}[Definable function]
Let $A\subseteq R^m$ and $B\subseteq R^n$. A function
\[
  f:A\to B
\]
is \emph{definable} if $A\in\mathcal S_m$, $B\in\mathcal S_n$, and its graph
\[
  \Graph(f)=\{(x,y)\in R^{m+n}:x\in A\text{ and }y=f(x)\}
\]
belongs to $\mathcal S_{m+n}$.
\end{definition}

In the Lean file, the function has type
\[
  \{x\in R^m:A(x)\}\to \{y\in R^n:B(y)\},
\]
so domain and codomain membership are built into the type of the function. The graph is the predicate
\[
  z\in\Graph(f)
  \quad\Longleftrightarrow\quad
  \exists h:A(z_L),\;z_R=f(z_L,h).
\]
Here $h$ is the proof that $z_L$ lies in the domain; it is needed in type theory because the input to $f$ is not merely a point of $R^m$, but a point together with a proof of membership in $A$.

\begin{lemma}[Definability of relational composition]
Let
\[
  G\subseteq R^{a+b},\qquad H\subseteq R^{b+c}
\]
be definable. Define
\[
  H\circ G
  =\{(x,z)\in R^{a+c}:\exists y\in R^b,
      (x,y)\in G\text{ and }(y,z)\in H\}.
\]
Then $H\circ G\in\mathcal S_{a+c}$.
\end{lemma}

\begin{proof}
Work in the larger space $R^{(a+c)+b}$ with coordinates written as $(x,z,y)$.
By reindexing coordinates, the set
\[
  G'=\bigl\{(x,z,y):(x,y)\in G\bigr\}
\]
is definable. Similarly,
\[
  H'=\bigl\{(x,z,y):(y,z)\in H\bigr\}
\]
is definable. Hence $G'\cap H'$ is definable. Project away the last block $y\in R^b$. The resulting subset of $R^{a+c}$ is exactly
\[
  \{(x,z):\exists y,
      (x,y)\in G\text{ and }(y,z)\in H\},
\]
which is the relational composition. Therefore it is definable.
\end{proof}

\begin{proposition}[Composition of definable functions]
Let $f:A\to B$ and $g:B\to C$ be definable functions. Then
\[
  g\circ f:A\to C
\]
is definable.
\end{proposition}

\begin{proof}
The domain $A$ and codomain $C$ are definable because they are respectively the domain of $f$ and codomain of $g$. For the graph, observe that
\[
  (x,z)\in\Graph(g\circ f)
\]
if and only if there exists $y\in R^n$ such that
\[
  (x,y)\in\Graph(f)
  \quad\text{and}\quad
  (y,z)\in\Graph(g).
\]
Thus
\[
  \Graph(g\circ f)=\Graph(g)\circ\Graph(f).
\]
The right hand side is definable by the previous lemma. Hence $g\circ f$ is definable.
\end{proof}

\section{Common definability devices}

For one-dimensional variables we write
\[
  x<y
\]
for the order relation on the corresponding coordinates.
For $n$-tuples and coordinate indices $i,j<n$, define:
\begin{align*}
  \dom_I(i) &= \{z\in R^n:z_i\in I\},\\
  \Graph_G(i,j) &= \{z\in R^n:(z_i,z_j)\in G\},\\
  \operatorname{Lt}(i,j) &= \{z\in R^n:z_i<z_j\},\\
  \operatorname{Eq}(i,j) &= \{z\in R^n:z_i=z_j\}.
\end{align*}
If $I$ and $G$ are definable, then all these sets are definable: $\dom_I(i)$ and $\Graph_G(i,j)$ are obtained by reindexing; $\operatorname{Lt}(i,j)$ is obtained by reindexing the definable order relation; and $\operatorname{Eq}(i,j)$ is a diagonal.

We also use the abbreviations
\begin{align*}
  \val(a,w,b) &= \{z:a<w<b\},\\
  \nhd(c,x,d) &= \{z:c<x<d\},\\
  \core_G(x,a,b,v) &= \{z:(x,v)\in G\text{ and }a<v<b\}.
\end{align*}
Each is definable whenever $G$ is definable, since it is built from finite intersections of definable sets.

The set $R^n$ itself is definable, because
\[
  R^n=R^n\setminus\emptyset.
\]
This is the content of the Lean theorem \texttt{setUniv\_mem}.

\section{Continuity points}

Let $I\subseteq R$ be definable and let $f:I\to R$ be definable. Let $G=\Graph(f)\subseteq R^2$.

\begin{definition}[Continuity at a point, interval form]
For $x\in R$, say that $f$ is continuous at $x$ on $I$ if $x\in I$ and for every interval $(a,b)$ containing $f(x)$ there is an interval $(c,d)$ containing $x$ such that
\[
  y\in I\text{ and }c<y<d
  \quad\Longrightarrow\quad
  a<f(y)<b.
\]
Equivalently, in graph language:
\[
\begin{aligned}
 x\in I\quad\text{and}\quad
 \forall a,b,v\;[&G(x,v)\land a<v<b\Rightarrow\\
 &\exists c,d\;(c<x<d\land
   \forall y,w\,(I(y)\land c<y<d\land G(y,w)\Rightarrow a<w<b))].
\end{aligned}
\]
This is the predicate \texttt{ContinuousAtOnGraph}.
\end{definition}

The Lean proof rewrites this universal statement using complements and existential projections. A proposed pair $(c,d)$ fails for fixed $(x,a,b,v)$ precisely when there are $y,w$ such that
\[
  y\in I,
  \qquad c<y<d,
  \qquad G(y,w),
  \qquad \neg(a<w<b).
\]
Define the counterexample set in $R^8$, with coordinates
\[
  (x,a,b,v,c,d,y,w),
\]
by
\[
\begin{aligned}
  C_8=\{(x,a,b,v,c,d,y,w):{}&y\in I,
  \;c<y<d,
  \;G(y,w),
  \;\neg(a<w<b)\}.
\end{aligned}
\]
This set is definable by finite intersections and complements.
Projecting away $(y,w)$ gives a definable set
\[
  C_6(x,a,b,v,c,d)
\]
expressing that $(c,d)$ is bad.
Thus
\[
  \neg C_6(x,a,b,v,c,d)
\]
expresses that $(c,d)$ works.

Now define
\[
\begin{aligned}
  G_6(x,a,b,v,c,d)\Longleftrightarrow{}&
  G(x,v)\land a<v<b\land c<x<d\land \neg C_6(x,a,b,v,c,d).
\end{aligned}
\]
This is definable. Projecting away $(c,d)$ gives a definable set
\[
  E_4(x,a,b,v)
\]
which says that there exists a working neighbourhood $(c,d)$ for the interval $(a,b)$ around the value $v$.

A bad continuity witness is a tuple $(x,a,b,v)$ such that
\[
  G(x,v)\land a<v<b\land \neg E_4(x,a,b,v).
\]
Call this definable set $B_4$. Projecting away $(a,b,v)$ gives a definable set $B_1\subseteq R$ of points where continuity fails. Therefore
\[
  \Cont(f)=I\cap (R\setminus B_1)
\]
is definable.

\begin{theorem}[Definability of the continuity locus]
Let $f:I\to R$ be a definable function on a definable set $I\subseteq R$. Then the set of points of $I$ at which $f$ is continuous is definable.
\end{theorem}

\begin{proof}
The preceding construction writes the discontinuity locus as a finite sequence of operations starting from the definable sets $I$, $G$, and $<$: finite intersections, complements, and existential projections. Each operation preserves definability in an o-minimal structure. Hence the bad set $B_1$ is definable. Its complement is definable, and intersecting with $I$ is definable. Therefore the continuity locus is definable.
\end{proof}

\section{Locally constant points}

\begin{definition}[Local constancy at a point]
The function $f$ is locally constant at $x$ on $I$ if $x\in I$ and there exists an interval $(c,d)$ containing $x$ such that for every $y\in I$,
\[
  c<y<d\quad\Longrightarrow\quad f(y)=f(x).
\]
In graph language this is:
\[
  x\in I\quad\text{and}\quad
  \exists c,d,v\;[G(x,v)\land c<x<d\land
  \forall y,w\,(I(y)\land c<y<d\land G(y,w)\Rightarrow w=v)].
\]
This is the predicate \texttt{LocallyConstantAtOnGraph}.
\end{definition}

For coordinates
\[
  (x,c,d,v,y,w)\in R^6,
\]
define the bad set
\[
\begin{aligned}
  C^{\mathrm{lc}}_6=\{(x,c,d,v,y,w):{}& y\in I,
  \;c<y<d,
  \;G(y,w),
  \;w\ne v\}.
\end{aligned}
\]
This is definable: the inequality $w\ne v$ is the complement of the diagonal $w=v$.
Projecting away $(y,w)$ gives a definable set
\[
  C^{\mathrm{lc}}_4(x,c,d,v)
\]
which says that a counterexample to constancy exists in the neighbourhood.
Hence
\[
  \neg C^{\mathrm{lc}}_4(x,c,d,v)
\]
says that the neighbourhood works.

Define the good set in $R^4$ by
\[
\begin{aligned}
  G^{\mathrm{lc}}_4(x,c,d,v)
  \Longleftrightarrow{}& G(x,v)\land c<x<d\land
  \neg C^{\mathrm{lc}}_4(x,c,d,v).
\end{aligned}
\]
This is definable. Projecting away $(c,d,v)$ gives a definable subset
\[
  G^{\mathrm{lc}}_1\subseteq R
\]
of points admitting such a neighbourhood and value. The locally constant locus is
\[
  \LC(f)=I\cap G^{\mathrm{lc}}_1.
\]

\begin{theorem}[Definability of the locally constant locus]
Let $f:I\to R$ be definable. Then the set of points of $I$ at which $f$ is locally constant is definable.
\end{theorem}

\begin{proof}
The set $C^{\mathrm{lc}}_6$ is built from $I$, $G$, the order relation, and the complement of a diagonal, using finite intersections. It is therefore definable. Projection gives $C^{\mathrm{lc}}_4$, complement gives the condition that no counterexample exists, another finite intersection gives $G^{\mathrm{lc}}_4$, and projection gives $G^{\mathrm{lc}}_1$. Finally $I\cap G^{\mathrm{lc}}_1$ is definable.
\end{proof}

\section{Locally monotone points}

We use the non-strict convention:
\[
  \text{locally increasing} = \text{locally nondecreasing},
  \qquad
  \text{locally decreasing} = \text{locally nonincreasing}.
\]

\begin{definition}[Local monotonicity at a point]
The function $f$ is locally increasing at $x$ if $x\in I$ and there exists an interval $(c,d)$ containing $x$ such that for all $y_0,y_1\in I$,
\[
  c<y_0<d,
  \quad c<y_1<d,
  \quad y_0<y_1
  \quad\Longrightarrow\quad
  f(y_0)\le f(y_1).
\]
Since we only have the strict order symbol, the conclusion is written as
\[
  \neg(f(y_1)<f(y_0)).
\]
In graph language, for values $w_0=f(y_0)$ and $w_1=f(y_1)$, the forbidden inequality is $w_1<w_0$.

Similarly, $f$ is locally decreasing at $x$ if the forbidden inequality is $w_0<w_1$.
These are the predicates \texttt{LocallyIncreasingAtOnGraph} and \texttt{LocallyDecreasingAtOnGraph}.
\end{definition}

Use coordinates
\[
  (x,c,d,y_0,y_1,w_0,w_1)
  \in R^7.
\]
First define the common base counterexample set
\[
\begin{aligned}
  C^{\mathrm{base}}_7=\{(x,c,d,y_0,y_1,w_0,w_1):{}&
  y_0\in I,
  \;y_1\in I,
  \;c<y_0<d,
  \;c<y_1<d,\\
  &G(y_0,w_0),
  \;G(y_1,w_1),
  \;y_0<y_1\}.
\end{aligned}
\]
This is definable by the closure properties.

The increasing-bad set is
\[
  C^{\mathrm{inc}}_7
  =C^{\mathrm{base}}_7\cap\{w_1<w_0\}.
\]
The decreasing-bad set is
\[
  C^{\mathrm{dec}}_7
  =C^{\mathrm{base}}_7\cap\{w_0<w_1\}.
\]
Both are definable.

Projecting away $(y_0,y_1,w_0,w_1)$ gives definable sets
\[
  C^{\mathrm{inc}}_3(x,c,d),
  \qquad
  C^{\mathrm{dec}}_3(x,c,d),
\]
which say that a bad pair exists in the neighbourhood $(c,d)$.
Thus
\[
  \neg C^{\mathrm{inc}}_3(x,c,d)
\]
says that $(c,d)$ witnesses local increasingness, and
\[
  \neg C^{\mathrm{dec}}_3(x,c,d)
\]
says that $(c,d)$ witnesses local decreasingness.

Define
\begin{align*}
  G^{\mathrm{inc}}_3(x,c,d)
  &\Longleftrightarrow c<x<d\land\neg C^{\mathrm{inc}}_3(x,c,d),\\
  G^{\mathrm{dec}}_3(x,c,d)
  &\Longleftrightarrow c<x<d\land\neg C^{\mathrm{dec}}_3(x,c,d).
\end{align*}
These are definable. Projecting away $(c,d)$ gives definable subsets
\[
  G^{\mathrm{inc}}_1,
  \qquad
  G^{\mathrm{dec}}_1
\]
of $R$. The final loci are
\[
  \LInc(f)=I\cap G^{\mathrm{inc}}_1,
  \qquad
  \LDec(f)=I\cap G^{\mathrm{dec}}_1.
\]

\begin{theorem}[Definability of local monotonicity loci]
Let $f:I\to R$ be definable. Then both sets
\[
  \{x\in I:f\text{ is locally increasing at }x\}
\]
and
\[
  \{x\in I:f\text{ is locally decreasing at }x\}
\]
are definable.
\end{theorem}

\begin{proof}
The common base counterexample set is obtained from $I$, $G$, and $<$ using only reindexing and finite intersections. Intersecting it with either $w_1<w_0$ or $w_0<w_1$ gives the increasing-bad and decreasing-bad sets. Each is definable. Projection gives the sets of triples $(x,c,d)$ for which a bad pair exists. Taking complements expresses that no bad pair exists. Intersecting with $c<x<d$ and then projecting away $(c,d)$ gives the points admitting a good neighbourhood. Finally intersecting with $I$ gives the desired loci. Each operation preserves definability, so both loci are definable.
\end{proof}

\section{Small sanity checks}

The Lean file ends with two simple definability checks:
\begin{enumerate}[label=(\roman*)]
  \item the empty set is definable in every arity, by the axiom $\emptyset\in\mathcal S_n$;
  \item the order relation $\{(x,y):x<y\}$ is definable, by the distinguished axiom \texttt{lt\_mem}.
\end{enumerate}

\section{Dictionary between the Lean file and this document}

\begin{center}
\begin{tabular}{ll}
\hline
Lean name & Mathematical object \\
\hline
\texttt{Power R n} & $R^n$ \\
\texttt{Power.append} & tuple concatenation $x\mathbin{{}^\frown}y$ \\
\texttt{Power.left}, \texttt{Power.right} & left and right coordinate blocks \\
\texttt{setInter}, \texttt{setUnion}, \texttt{setCompl} & $\cap$, $\cup$, complement \\
\texttt{DenseLinearOrderNoEndpoints} & dense linear order without endpoints \\
\texttt{OMinimalStructure} & family $(\mathcal S_n)_n$ with closure axioms \\
\texttt{FunctionGraph} & graph of a definable function \\
\texttt{relation\_comp\_mem} & relational composition preserves definability \\
\texttt{DefinableFunction.comp} & composition of definable functions \\
\texttt{ContinuousPoints} & continuity locus \\
\texttt{LocallyConstantPoints} & local constancy locus \\
\texttt{LocallyIncreasingPoints} & locally nondecreasing locus \\
\texttt{LocallyDecreasingPoints} & locally nonincreasing locus \\
\hline
\end{tabular}
\end{center}

\end{document}
